% Module 4 section file. Included by ../module4_alfven_continuum_eigenmodes.tex.
\section{Exercises and computational laboratories}
\label{sec:exercises}

\subsection{Conceptual exercises}
\begin{enumerate}[leftmargin=1.8em]
\item Explain in words why a transverse displacement that is constant along a uniform field produces no shear-Alfv\'en restoring force.
\item Distinguish magnetic pressure from magnetic tension using Eq.~\eqref{eq:magnetic_pressure_tension}.
\item Explain why a fluid-MHD eigenmode is a legitimate component of a kinetic-MHD model.
\item Explain why $\omega_A(s)$ forms a continuum when $s$ is continuous.
\item Distinguish a rational surface, a continuum-resonant surface, and a neighboring-harmonic crossing surface.
\item Explain why a local avoided crossing does not prove the existence of a global gap mode.
\item Explain why mode amplitude is undetermined in a homogeneous linear eigenproblem.
\item Explain the difference between verification and validation using the current benchmark as an example.
\end{enumerate}

\subsection{Derivation exercises}
\begin{enumerate}[leftmargin=1.8em]
\item Starting from Eqs.~\eqref{eq:mhd_induction} and \eqref{eq:v_xi}, derive Eq.~\eqref{eq:b1_xi}.
\item Starting from the adiabatic pressure equation, derive Eq.~\eqref{eq:p1_xi}.
\item Repeat the uniform shear-Alfv\'en derivation for displacement in the $y$ direction.
\item Show explicitly that the uniform shear-Alfv\'en wave has equal cycle-averaged kinetic and magnetic energy.
\item Derive Eq.~\eqref{eq:kparallel_benchmark} from the stated Fourier convention and field-line pitch, carefully documenting the sign convention.
\item Derive the logarithmic continuum slope in Eq.~\eqref{eq:continuum_log_slope}.
\item Diagonalize Eq.~\eqref{eq:generic_two_mode_matrix} and derive both eigenvalues and the mixing-angle relation.
\item Derive the crossing transform and crossing frequency in Eqs.~\eqref{eq:tae_crossing_iota}--\eqref{eq:tae_gap_center}.
\item Derive the Euler--Lagrange equations from Eqs.~\eqref{eq:reduced_energy}--\eqref{eq:reduced_norm}.
\item Prove the discrete positivity identity in Eq.~\eqref{eq:discrete_K_positive}.
\end{enumerate}

\subsection{Numerical exercises}
\begin{enumerate}[leftmargin=1.8em]
\item Reproduce the baseline density, transform, Alfv\'en speed, and uncoupled continuum arrays directly from the JSON file without calling the project solver.
\item Verify the crossing location numerically and compare the exact $s_c$ with the nearest cell center for $N=64,128,256,512$.
\item Set $\epsilon_c=0$ and verify that the two block families can be solved independently.
\item Scan $B_0$ and confirm that every frequency is linear in $B_0$ while normalized eigenvectors remain unchanged to numerical precision.
\item Scan a uniform density multiplier and confirm the inverse-square-root scaling.
\item Compute the observed convergence order for modes 1--4 using at least four grids and Richardson extrapolation.
\item Calculate the discrete orthogonality matrix $h\mathbf X^{\mathsf T}\mathbf X$.
\item Reproduce Table~\ref{tab:harmonic_diagnostics} from the eigenmode CSV.
\item Track modes through a coupling scan using maximum eigenvector overlap. Compare with sorted-index tracking.
\item Replace the arithmetic face average by a harmonic average and determine the convergence difference for the smooth baseline coefficient.
\end{enumerate}

\subsection{Model-development laboratories}
\begin{enumerate}[leftmargin=1.8em]
\item Tune $\epsilon_c$, $w_c$, and $\epsilon_r$ to produce a mode whose frequency lies in the local gap and whose radial density peaks near $s_c$. Preserve the original case and create a separately versioned case.
\item Generalize the operator to $M$ harmonics with a sparse coupling graph. Verify symmetry and the decoupled limit for every edge of the graph.
\item Add a nonuniform mass matrix $\mathbf M$ and solve $\mathbf A\mathbf x=\omega^2\mathbf M\mathbf x$. Define the correct normalization and residual.
\item Construct a cylindrical analytic benchmark with constant $D$ and constant potential so that mixed-boundary eigenvalues are known exactly.
\item Implement a nonuniform radial grid and derive face distances and cell widths carefully.
\item Create an eigenvalue PINN using the exact boundary-satisfying trial function from Worked Example 10. Compare strong-form and variational losses.
\item Train multiple neural modes using orthogonality constraints. Test whether the network repeatedly collapses to the lowest mode.
\item Import an equilibrium-derived $\iota(s)$ and density profile while retaining the surrogate coupling. Quantify which changes are simple rescalings and which alter mode structure.
\item Replace the surrogate coupling with coefficients derived from a controlled large-aspect-ratio toroidal model.
\item Design a validation case against a trusted external eigenmode code, including equilibrium conversion, harmonic convention, normalization, and uncertainty accounting.
\end{enumerate}

\section{Chapter summary}
\label{sec:summary}

\begin{physicsbox}
\textbf{Physical chain.}
\begin{equation*}
\begin{aligned}
\text{magnetic tension}
&\rightarrow \omega=|k_\parallel|v_A
\rightarrow \omega_m(s)\\
&\rightarrow \text{Alfv\'en continuum}
\rightarrow \text{geometric harmonic coupling}\\
&\rightarrow \text{spectral gaps}
\rightarrow \text{global eigenmodes}.
\end{aligned}
\end{equation*}
\end{physicsbox}

The shear-Alfv\'en wave is a transverse fluid-MHD oscillation restored mainly by magnetic tension. In a uniform medium it satisfies
\begin{equation}
\omega^2=k_\parallel^2v_A^2,
\qquad
v_A=\frac{B_0}{\sqrt{\mu_0\rho_0}}.
\end{equation}
In a toroidal plasma, field-line pitch and density vary across magnetic surfaces. For the benchmark convention,
\begin{equation}
k_{\parallel,m}=\frac{n-m\iota}{R_0},
\end{equation}
so each harmonic has a radially varying local continuum branch.

Toroidal and three-dimensional geometry couples Fourier harmonics. A symmetric two-mode matrix produces avoided crossing, harmonic mixing, and a local gap. Neighboring harmonics $m$ and $m+1$ cross at
\begin{equation}
\iota_c=\frac{2n}{2m+1}.
\end{equation}
A local gap alone does not establish a global eigenmode. Radial derivatives and boundary conditions must be solved.

The project's executable benchmark uses two coupled radial envelopes:
\begin{equation}
-\frac{d}{ds}\left(D\frac{d\boldsymbol\xi}{ds}\right)
+\mathbf H_{\mathrm{loc}}\boldsymbol\xi
=\omega^2\boldsymbol\xi.
\end{equation}
It follows from a positive quadratic functional and is self-adjoint under zero axis flux and zero edge value. A cell-centered finite-volume discretization produces a real symmetric sparse matrix.

The baseline solver is strongly verified as a reduced numerical benchmark: it passes symmetry, decoupling, exact Alfv\'en scaling, residual, regression, and grid-refinement tests; Python reference tests pass, and MATLAB parity is defined for environments with MATLAB. The frozen output demonstrates coupled branches and convergent mixed-harmonic eigenmodes.

The baseline should not be overinterpreted. Mode 2 lies within the local crossing-cell interval but is not strongly localized at the crossing and is dominated by one harmonic. The current result is therefore not a validated predictive stellarator TAE. It is a transparent conventional target for the next stages: many-harmonic geometry, equilibrium-derived coefficients, kinetic drive and damping, and eigenvalue PINNs.

\section{Minimum interview-level explanation}
A concise but technically accurate explanation is:

\begin{overviewbox}
A shear-Alfv\'en wave bends magnetic field lines, and magnetic tension restores them. Its local frequency is approximately $\omega_A=|k_\parallel|v_A$. In a toroidal plasma, rotational transform, magnetic geometry, and density vary across flux surfaces, so the local Alfv\'en frequency forms a continuum. Toroidal or stellarator geometry couples Fourier harmonics and opens gaps at avoided crossings. Global Alfv\'en eigenmodes can form through radial coupling and boundary conditions, sometimes inside those gaps. Module 4 solves a verified reduced fluid-MHD eigenproblem for that structure. Module 5 then determines whether energetic particles drive or damp the mode through kinetic resonances.
\end{overviewbox}

\section{Completion checklist for the reader}
A reader is prepared to move to Module 5 when they can answer all of the following without relying on memorized phrases:
\begin{itemize}[leftmargin=1.6em]
\item What physical force restores a shear-Alfv\'en wave?
\item Why is $v_A$ proportional to $B$ and inversely proportional to $\sqrt\rho$?
\item How is $k_\parallel$ obtained from $(m,n,\iota)$ under a stated convention?
\item Why does radial nonuniformity create a continuum?
\item What is the difference between a local continuum branch and a global eigenmode?
\item How does a $2\times2$ coupling matrix open a gap?
\item Why is a mode frequency insufficient for classifying a TAE?
\item Why is the reduced operator self-adjoint?
\item How are the axis and edge boundary conditions discretized?
\item What do residual, grid convergence, and code parity each establish?
\item Which information must Module 4 pass to a kinetic resonance calculation?
\end{itemize}
